\documentclass[12pt]{lexlean}
\usepackage{prism.arch}

\title{UorNotary.Product --- Product Identity, Free Commons Scope \& Claim Boundaries}
\author{UOR Foundation Commons Engineering}
\lexmodule{UorNotary.Product}
\lexlanguage{1.2}

\begin{document}
\maketitle

\section{Product Identity \& Zero-Lock-In Free Commons Mandate}

\begin{lean}
namespace UorNotary.Product

def crateName : String := "prism_mcp_notary"
def packageName : String := "uor-mcp-notary"
def modulePrefix : String := "UorNotary"
def languageVersion : String := "1.2"
def mcpProtocolVersion : String := "2026-07-28"
def uorAddrVersion : String := "0.2.0"
def spdxLicense : String := "MIT OR Apache-2.0"

inductive GovernanceModel where
  | OpenFoundationCommons
deriving DecidableEq, Repr

def activeGovernanceModel : GovernanceModel :=
  GovernanceModel.OpenFoundationCommons

inductive ClaimDomain where
  | ImplementationAcceptance
  | PayloadBlindness
  | UnicodeSmugglingRefusal
  | OfflineVerification
  | ZeroLockInFreeCommons
deriving DecidableEq, Repr

structure ClaimBoundary where
  domain          : ClaimDomain
  formalGuarantee : String
  explicitNonScope : String
  registerBinding : String
deriving DecidableEq, Repr

def boundaryFor (d : ClaimDomain) : ClaimBoundary :=
  match d with
  | .ImplementationAcceptance =>
    { domain := .ImplementationAcceptance
    , formalGuarantee :=
        "just vv passes cleanly (lexlean fmt --check, lexlean check, lexlean verify, prismpm check, prismpm build, golden vectors, planted defects)."
    , explicitNonScope :=
        "Does NOT prove host OS kernel integrity, hardware KMS physical security, or remote MCP server internal security."
    , registerBinding := "AU-13"
    }
  | .PayloadBlindness =>
    { domain := .PayloadBlindness
    , formalGuarantee :=
        "Neither HostEffect (E1-E7) nor Observability can carry raw MCP text, arguments, or outputs (noRawPayloadInHostEffects, redactedHasOneValue)."
    , explicitNonScope :=
        "Does NOT protect against a user deliberately modifying the thin host adapter to log raw inputs before calling core-wasm."
    , registerBinding := "EF-05"
    }
  | .UnicodeSmugglingRefusal =>
    { domain := .UnicodeSmugglingRefusal
    , formalGuarantee :=
        "admitSurface deterministically refuses all forbidden scalar classes before canonicalization under StrictV1."
    , explicitNonScope :=
        "Does NOT perform semantic LLM prompt-injection detection on ordinary visible ASCII/Unicode prose (Risk R8 Residual)."
    , registerBinding := "UR-01"
    }
  | .OfflineVerification =>
    { domain := .OfflineVerification
    , formalGuarantee :=
        "Any third party can verify DSSE PAE signatures, recompute kappa-labels via uor-addr, and check .holo / Merkle proofs with zero network calls (TC-05/TC-06)."
    , explicitNonScope :=
        "Does NOT require any central paid server; commercial operators only compete on uptime, HSMs, and enterprise alert routing."
    , registerBinding := "WT-04"
    }
  | .ZeroLockInFreeCommons =>
    { domain := .ZeroLockInFreeCommons
    , formalGuarantee :=
        "Every protocol rule, CLI sidecar, self-hostable notary, offline verifier, and public Watchtower runs free under MIT OR Apache-2.0."
    , explicitNonScope :=
        "Paid hosted compliance tiers are strictly out-of-scope of this repository."
    , registerBinding := "WT-06"
    }

def isBoundaryWellFormed (b : ClaimBoundary) : Bool :=
  (!b.formalGuarantee.isEmpty) &&
  (!b.explicitNonScope.isEmpty) &&
  (!b.registerBinding.isEmpty)

theorem claimBoundaryComplete (d : ClaimDomain) :
    isBoundaryWellFormed (boundaryFor d) = true := by
  cases d <;> rfl

theorem freeCommonsZeroLockIn :
    activeGovernanceModel = GovernanceModel.OpenFoundationCommons ∧
    spdxLicense = "MIT OR Apache-2.0" := by
  exact ⟨rfl, rfl⟩

theorem crateAndPrefixDiscipline :
    crateName = "prism_mcp_notary" ∧ modulePrefix = "UorNotary" := by
  exact ⟨rfl, rfl⟩

theorem boundaryDomainPreserved (d : ClaimDomain) :
    (boundaryFor d).domain = d := by
  cases d <;> rfl

end UorNotary.Product
\end{lean}

\end{document}
